% !TeX encoding = UTF-8
% !TeX TS-program = latexmk
% This is file `main-en.tex'.
% !Read the user manual before use. Copyright (C) 2021-2026 by Linrong Wu.

% ================ %
%      Preamble    %
% ================ %
\PassOptionsToPackage{cmyk}{xcolor}% Avoid color drift when generating PDFs with CMYK colors.
\documentclass[hyperref, UTF8, CJK, aspectratio=169]{beamer}

% --------BIT Beamer theme package--------
% ----------------
\usetheme[
	ColorDisplay=BITB, % BITA ⚙️ | BITB | Custom → theme color display
	% BlockDisplay=colorful, % colorful ⚙️ | followtheme | allgrey → block color display
	CodeTheme=listing, % listing ⚙️ | minted | minted2 → code highlighting engine
	% MintedStyle=dracula, % lightmode ⚙️ | darkmode | ⟨custom⟩ → minted style; active only with CodeTheme=minted or minted2
	LanguageMode=en, % cn ⚙️ | en → language mode
	Miniframes=follow, % follow ⚙️ | separate | none → headline mini-frame behavior
	% NavigationTool=1-2-3, % 1-2-3 ⚙️ | ⟨see manual⟩ | none → footline navigation tools
	% FontTheme=Auto, % Auto ⚙️ | Ubuntu | Win | Mac | Fandol | Source-Han | ... | Custom → font theme
	% MathFont=LM, % LM ⚙️ | XITS | ⟨custom⟩ → math font
	% BIBMode=none, % biber ⚙️ | none → bibliography engine
	% BIBStyle=numeric-comp, % numeric-comp ⚙️ | ⟨custom⟩ → bibliography style; inactive when BIBMode=none
	% ContentMuticols=true, % true ⚙️ | false → two-column table of contents
	Background=BIT-Lite, % BIT-Full ⚙️ | BIT-Lite | Custom | none → background display
]{bit}

% --------Package imports--------
% ----------------
\usepackage{array}
\usepackage{algorithm,algorithmic}
\usepackage[english]{babel}
\usepackage{color}      % color content
\usepackage{url}        % hyperlinks
\usepackage{multicol,multirow}
\usepackage{ulem} % ulem: underline and strikeout helpers.
\usepackage{booktabs}
\usepackage{cprotect}
\usepackage{makecell}
\usepackage{listings}
\usepackage{subcaption}
\usepackage{varioref,cleveref}
%\usepackage[active,tightpage]{preview}
%\PreviewEnvironment{bitcode}

% --------Custom commands--------
% Define frequently used commands here.
% ----------------
\newcommand{\fverb}[1]{\texttt{#1}}

% --------Title page metadata--------
% [<in footline>] is shown in the footline; {<in title page>} is shown on the cover.
% ----------------
\title[A brief English demo for BIT Beamer Theme]{A brief English demo}
\subtitle{BIT Beamer Theme}
\author[BIT Beamer Theme]{BIT Beamer Theme}
\institute{%
	School of Management and Economics\\
	Beijing Institute of Technology\\
	\mailbit{lr.wu.interact@outlook.com}
}
\date{\today}

% ---------------- %
%      Content     %
% ---------------- %
\begin{document}

% --------Table of contents--------
% Optional.
% ----------------
%	\begin{frame}{Contents}
%		%\transfade
%		\tableofcontents
%	\end{frame}

% --------Section: Introduction--------
% ----------------
\section{Introduction}
\subsection{The Project}
\begin{frame}{Info.}
	\faShare*\enspace{\color{bitb}lr.wu.interact@outlook.com}
	\faGithub\enspace{\color{bitb}\url{https://github.com/FvNCCR228/bit_beamer_theme}}

	\vspace{2ex}
	\begin{itemize}
		\item A compact English demo for BIT Beamer Theme.
		\item Built with XeLaTeX and the Beamer theme system.
		\item Derived from SCU Beamer Theme with BIT visual identity.
	\end{itemize}
\end{frame}

% --------Section: Block examples--------
% ----------------
\section{Blocks}
\subsection{Math Blocks}
\begin{frame}[fragile,allowframebreaks]{Math Blocks}
	\begin{bittheorem}{A Theorem}[theorem1]
		\begin{equation}
		\dfrac{1}{n} \sum_{k=1}^{n} X_k - \dfrac{1}{n} \sum_{k=1}^{n} E(X_k) \stackrel{\;P\;}{\longrightarrow} 0
		\end{equation}
	\end{bittheorem}
	\begin{bitproof}{}
		A proof block.
	\end{bitproof}
	\begin{bitexample}{An Example}[example1]
		An example block.
	\end{bitexample}
	\begin{bitalgorithm}{An Algorithm}[algorithm1]
		\begin{algorithmic}[1]
			\REQUIRE \LaTeX{}
			\ENSURE Computer
			\STATE ST
			\STATE A
			\STATE TE
			\RETURN Beamer
		\end{algorithmic}
	\end{bitalgorithm}
	\begin{bitdefinition}{A Definition}[definition1]
		A definition block.
	\end{bitdefinition}
	\begin{bitaxiom}{An Axiom}[axiom1]
		An axiom block. Reference to~\Vref{def:definition1}
	\end{bitaxiom}
	\begin{bitproperty}{A Property}[property1]
		A property block. Reference to~\Cref{axio:axiom1}
	\end{bitproperty}
	\begin{bitproposition}{A Proposition}[proposition1]
		A proposition block. Reference to~\vref{prope:property1}
		\begin{equation}
		\Delta x \Delta p \geq \dfrac{h}{4\pi}
		\end{equation}
		where $h$ is Planck's constant.
	\end{bitproposition}
	\begin{bitlemma}{A Lemma}[lemma1]
		A lemma block. Reference to~\cref{propo:proposition1}
	\end{bitlemma}
	\begin{bitcorollary}{A Corollary}[corollary1]
		A corollary block.
	\end{bitcorollary}
	\begin{bitremark}{}[remark1]
		A remark block.
	\end{bitremark}
	\begin{bitcondition}{A Condition}[condition1]
		A condition block.
	\end{bitcondition}
	\begin{bitconclusion}{A Conclusion}[conclusion1]
		A conclusion block.
	\end{bitconclusion}
	\begin{bitassumption}{An Assumption}[assumption]
		An assumption block.
	\end{bitassumption}
\end{frame}

\begin{frame}{A Starred Block}
\begin{bittheorem}*<1-4>[after title=(after title: Theorem)]{A Starred Theorem Block}[staredblock1]
	\begin{itemize}[<+->]
		\item One
		\item Two~\only<2>{Two}
		\item \alert<3>{Three}
		\item Four
	\end{itemize}
\end{bittheorem}
\begin{bittheorem}<2-5>[after title=(after title: Theorem)]{Another Starred Theorem Block}*[staredblock2]
	\begin{itemize}
		\item Five
		\item Six~\only<3>{Six}
		\item \alert<3>{Seven}
		\item Eight
	\end{itemize}
\end{bittheorem}
\end{frame}

\subsection{Source Code Block}
\begin{frame}[fragile]{Source Code Block}{With frame option ``fragile''}
	\onslide<2>
	\begin{bitcode}{A Cpp Program.}[cppcode]{c}
		#include <iostream>
		int main()
		{
			std::cout << "Hello World!" << std::endl;
			std::cin.get();
		}
	\end{bitcode}
	\onslide<1>
	\begin{bitcode}{A Python Program.}[pythoncode]{python}
		for i in range(1,5):
			for j in range(1,5):
				for k in range(1,5):
					if( i != k ) and (i != j) and (j != k):
						print (i,j,k)
	\end{bitcode}
\end{frame}

\begin{frame}[fragile]{A Starred Source Code Block}
	\begin{bitcode*}{A Starred Block.}[starredcode]{c}
		#include <iostream>
		int main()
		{
			std::cout << "Hello World!" << std::endl;
			std::cin.get();
		}
	\end{bitcode*}
	\begin{bitcode}{Another Starred Theorem Block.}*[starredpythoncode]{python}
		for i in range(1,5):
			for j in range(1,5):
				for k in range(1,5):
					if( i != k ) and (i != j) and (j != k):
						print (i,j,k)
	\end{bitcode}
\end{frame}

\begin{frame}[fragile]{Line Emphasis}
	\begin{bitcode}{Emphasized C++ lines.}[highlightcpp]{c}[emph={main,cin},emphstyle=\color{PrimaryC}\bfseries]
		#include <iostream>
		int main()
		{
			std::cout << "Hello World!" << std::endl;
			std::cin.get();
		}
	\end{bitcode}
	\begin{bitcode}{Emphasized Python lines.}[highlightpython]{python}[emph={range,print},emphstyle=\color{PrimaryC}\bfseries]
		for i in range(1,5):
			for j in range(1,5):
				for k in range(1,5):
					if( i != k ) and (i != j) and (j != k):
						print(i,j,k)
	\end{bitcode}
	References to~\cref{code:highlightcpp} and~\cref{code:highlightpython}.
\end{frame}

\begin{frame}[fragile]{\LaTeX{} Comment}{Escapeinline}
	{\footnotesize Use native comments for ordinary notes; use escape delimiters only for inline \LaTeX{}.}
	\vspace{-.5ex}
	\begin{bitcode}{Comment.}[commentcpp]{c}
		#include <iostream>
		int main()
		{// $\pi$
			std::cout << "Hello World!" << std::endl; |\# \textsf{\LaTeX{} out hEllo wOrld}|
			|\colorbox{bitc!60}{$\sum_\pi^\phi \alpha + \Gamma$}| std::cin.get(); 
		}
	\end{bitcode}
	\begin{bitcode}{Comment}[commentpython]{python}'@@'
		for i in range(1,5):
			for j in range(1,5): @$\sum_\pi^\phi \alpha + \Gamma$@
				for k in range(1,5): # $\sum_\pi^\phi \alpha + \Gamma$
					if( i != k ) and (i != j) and (j != k):
						print (i,j,k)
	\end{bitcode}
\end{frame}

\begin{frame}[fragile]{Overlay \& Label}{Escapeinline}
	\begin{bitcode}{Comment.}[overlaycpp]{c}
		#include <iostream>
		int main()
		{
			std::cout << "Hello World!" << std::endl; // \only<1>{Value 1}\only<2>{Value 2}
			std::cin.get();
		}
	\end{bitcode}
	\begin{bitcode}{Comment}[overlaypython]{python}[emph={if},emphstyle=\color{PrimaryC}\bfseries]'@@'
		for i in range(1,5):
			for j in range(1,5): 
				for k in range(1,5):
					if( i != k ) and (i != j) and (j != k): @\label{line:qg}@
						print (i,j,k)
	\end{bitcode}
  Reference to Line~\ref{line:qg}, the if statement.
\end{frame}

\begin{frame}[fragile]{Source Code From File}
	\bitcodeinput{Source Code From File}{c}{A cpp.cpp}
\end{frame}

% --------Section: References--------
% ----------------
% \section{Reference}
% \subsection{Reference}
% \begin{frame}[allowframebreaks]{Reference}
% 	\nocite{*}
% 	\printbibliography[heading=none]
% \end{frame}

% --------Section: Acknowledgement--------
% ----------------
\section{Acknowledgement}
\subsection{Acknowledgement}
\begin{frame}{Acknowledgement}
	\begin{center}
		BIT Beamer Theme is derived from SCU Beamer Theme, and this sample keeps \texttt{SCU-Version} as source-lineage attribution.\\[1ex]
		Thanks to the authors and maintainers of Beamer, Tcolorbox, BibLaTeX, FontAwesome, and related packages.\\[1ex]
		Thanks to the community discussions that helped shape the original template behavior.\\[1ex]
	\end{center}
\end{frame}

\begin{frame}
	\centering
	\Huge Thanks!
\end{frame}

\end{document}

%% End of file `main-en.tex'.
